% گزارش پروژه Patchcheck — امنیت داده و شبکه
% برای کامپایل: xelatex report.tex (به قالب styles/hw و پوشه‌های fonts/ و images/ از تمپلیت شریف نیاز دارد)
\documentclass[12pt]{exam}

\usepackage{styles/hw}

% ---------- سربرگ اختصاصی ----------
\renewcommand{\heading}{%
	\begin{center}
		\textbf{\Huge به نام خدا}

		\vspace{0.2cm}

		\includegraphics[width=0.4\textwidth]{images/sharif.png}

		\vspace{0.2cm}

		\textbf{\Huge امنیت داده و شبکه}

		\vspace{0.25cm}

		\textbf{\Large گزارش پروژهٔ اول: اعتبارسنجی خودکار وصله‌های امنیتی و آزمون رگرسیون در \lr{CI/CD}}

		\vspace{0.2cm}

		\large\textbf{دانشکده مهندسی کامپیوتر}

		\vspace{0.1cm}

		\large دانشگاه صنعتی شریف

		\vspace{0.2cm}

		\noindent\rule{\linewidth}{1pt}
	\end{center}

	مدرس:

	\textbf{سولماز سلیمی}

	\vspace{0.2cm}

	اعضای گروه:

	\textbf{امیرحسین شایان \lr{402170981} - امیرعلی بیگی \lr{402110554}}
}

\begin{document}

\begin{titlepage}
	\begin{center}
		\heading{}
	\end{center}
\end{titlepage}

\newpage
\newpage


% ===================== چکیده =====================
\section{چکیده}

این پروژه ابزاری خودکار به نام \codeword|Patchcheck| می‌سازد که در پایپ‌لاین \lr{CI/CD}
یک وصلهٔ امنیتی را \textbf{واقعاً اجرا} می‌کند و به دو پرسش صورت‌پروژه پاسخ می‌دهد:
(۱) آیا خطای امنیتی قبلی بعد از وصله همچنان رخ می‌دهد؟ و (۲) آیا وصله، رفتار درست
برنامه را روی ورودی‌های معتبر حفظ کرده است؟ ابزار نسخهٔ قبل و بعد از وصله را جداگانه
به‌صورت \lr{ELF} روی لینوکس \lr{x86\_64} و در دو حالت \lr{Release} و \lr{Sanitized}
می‌سازد، ورودی‌های معتبر/مرزی/نامعتبر/خطاساز را در یک کانتینر ایزوله با محدودیت منابع
اجرا می‌کند، خروجی را با یک مرجع مستقل (\lr{oracle}) مقایسه می‌کند، خطاها را طبقه‌بندی
می‌کند و یک کد خروج استاندارد تولید می‌کند. ابزار روی پنج برنامهٔ مستقل آزمایش شده
است؛ سه هدف واقعی (\codeword|tlvparse|، \codeword|giflib|، \codeword|libpng|) و دو
نمونهٔ بی‌خطر برای خودآزمایی چارچوب (\codeword|bytecount|، \codeword|lineparse|).
هر هدف در یک مخزن مستقل گیت‌هاب قرار دارد و با \lr{Pull Request} ارزیابی می‌شود.

% ===================== مخازن =====================
\section{مخازن پروژه}

پروژه از یک مخزن اصلی (خود ابزار) و پنج مخزن مستقل هدف تشکیل شده است:

\begin{center}
\begin{tabular}{lll}
\toprule
نام & نقش & آدرس \\
\midrule
\codeword|patchcheck| & ابزار اصلی + همهٔ هدف‌ها و نسخه‌ها & \url{https://github.com/amirhoseinshayan/patchcheck} \\
\codeword|patchcheck-tlvparse| & هدف خودِ گروه & \url{https://github.com/amirhoseinshayan/patchcheck-tlvparse} \\
\codeword|patchcheck-giflib| & پروژهٔ واقعی & \url{https://github.com/amirhoseinshayan/patchcheck-giflib} \\
\codeword|patchcheck-libpng| & کتابخانهٔ واقعی & \url{https://github.com/amirhoseinshayan/patchcheck-libpng} \\
\codeword|patchcheck-bytecount| & نمونهٔ بی‌خطر (فیچر) & \url{https://github.com/amirhoseinshayan/patchcheck-bytecount} \\
\codeword|patchcheck-lineparse| & نمونهٔ بی‌خطر (فیچر) & \url{https://github.com/amirhoseinshayan/patchcheck-lineparse} \\
\bottomrule
\end{tabular}
\end{center}

در مخزن اصلی، همهٔ هدف‌ها زیر پوشهٔ \codeword|targets/| با نسخه‌های دقیق‌شان نگه‌داری
می‌شوند. در مخازن مستقل، کد همان هدف روی شاخهٔ \codeword|main| (نسخهٔ پایه) است و هر
وصله به‌صورت یک \lr{Pull Request} جداگانه ارزیابی می‌شود.

% ===================== معماری =====================
\section{معماری کلی ابزار}

ابزار یک برنامهٔ پایتون است که برای هر ارزیابی این مراحل را طی می‌کند:

\begin{enumerate}
\item \textbf{برداشت نسخه‌ها}: دو نسخهٔ \codeword|base| (قبل از وصله) و
\codeword|candidate| (بعد از وصله) را از مخزن گیت با همان \lr{commit} مشخص‌شده استخراج
و \lr{diff} بین آن‌ها را ذخیره می‌کند.
\item \textbf{ساخت}: هر دو نسخه در دو حالت \codeword|release| و \codeword|sanitized|
(با \lr{AddressSanitizer}/\lr{UBSan}) داخل یک کانتینر لینوکس \lr{x86\_64} کامپایل
می‌شوند.
\item \textbf{اجرا}: هر ورودی آزمون روی هر چهار بیلد (۲ نسخه × ۲ حالت) در یک کانتینر
ایزوله با محدودیت زمان/حافظه/خروجی اجرا می‌شود و \lr{stdout}، \lr{stderr}، کد خروج و
سیگنال ثبت می‌شود.
\item \textbf{تحلیل}: خروجی با مرجع مستقل مقایسه و خطاها بر اساس شواهد واقعی طبقه‌بندی
می‌شوند.
\item \textbf{تصمیم}: بر اساس همهٔ شواهد، یکی از کدهای خروج استاندارد تولید می‌شود.
\item \textbf{پیگیری}: یافته‌های تأییدشده حداقل سه بار بازپخش می‌شوند و بستهٔ بازتولید
ساخته می‌شود.
\end{enumerate}

نکتهٔ مهم معماری این است که ابزار به نام تابع یا ساختار یک برنامهٔ خاص وابسته نیست؛
افزودن هدف جدید فقط با یک فایل پیکربندی \lr{JSON} انجام می‌شود.

% ===================== الزامات پایه =====================
\section{دامنهٔ پشتیبانی و الزامات پایه}

\subsection{هدف از نوع C/C++ و خروجی ELF روی لینوکس}

صورت‌پروژه می‌گوید هدف باید برنامهٔ کامپایل‌شدهٔ \lr{C}/\lr{C++} با خروجی \lr{ELF}
روی لینوکس \lr{x86\_64} باشد. این شرط نه فقط در مستندات، بلکه در خود کد اعمال شده
است؛ تابع \codeword|check\_elf| بعد از ساخت، سرآیند باینری را می‌خواند و اگر فرمت
\lr{ELF} ۶۴ بیتی یا معماری \lr{x86\_64} نباشد، ساخت را رد می‌کند:

\begin{latin}
\begin{lstlisting}[language=Python]
def check_elf(binary):
    header = Path(binary).read_bytes()[:64]
    if len(header) < 64 or header[:4] != b'\x7fELF' or header[4:6] != b'\x02\x01':
        raise ConfigError('build output must be a 64-bit little-endian ELF executable')
    kind, machine = struct.unpack_from('<HH', header, 16)
    if machine != 62 or kind not in (2, 3):   # 62 = x86_64
        raise ConfigError(f'ELF must be x86_64 EXEC/DYN, got machine={machine}')
\end{lstlisting}
\end{latin}

همچنین قبل از ساخت، محیط داخل کانتینر بررسی می‌شود که واقعاً \lr{Linux} و
\lr{x86\_64} است.

\subsection{ورودی از فایل، آرگومان یا ورودی استاندارد}

ابزار هر سه حالت ورودی را پشتیبانی می‌کند. حالت با \codeword|input.mode| در پیکربندی
مشخص می‌شود و هر سه حالت پوشش داده شده‌اند:

\begin{center}
\begin{tabular}{lll}
\toprule
هدف & حالت ورودی & نمونه \\
\midrule
\codeword|tlvparse|، \codeword|giflib|، \codeword|libpng| & فایل & \codeword|\{binary\} \{input\}| \\
\codeword|bytecount| & ورودی استاندارد & \codeword|\{binary\} stdin| \\
\codeword|lineparse| & آرگومان & \codeword|\{binary\} argv \{input\}| \\
\bottomrule
\end{tabular}
\end{center}

\subsection{قابلیت تعمیم فقط با پیکربندی}

کل تعریف هر هدف — دستور ساخت، مسیر باینری، نحوهٔ ارسال ورودی و مجموعهٔ آزمون‌ها — در
یک فایل \lr{JSON} است. برای نمونه بخش ساخت و ورودی هدف \codeword|libpng|:

\begin{latin}
\begin{lstlisting}[language=Python]
"builds": { "release": { "compiler": "gcc", "flags": ["-O2", "-D_FORTIFY_SOURCE=2", ...],
            "command": ["{compiler}", "{flags}", "-Isrc", "src/png.c", ..., "-o", "pnginfo", "-lz", "-lm"] } },
"input": { "mode": "file", "argv": ["{binary}", "{input}"] }
\end{lstlisting}
\end{latin}

% ===================== آزمون پویا =====================
\section{آزمون پویا و گزارش دقیق خطا}

صورت‌پروژه تأکید می‌کند که «اجرای واقعی» اجباری است و اکتفا به تحلیل ایستا پذیرفته
نیست. این بخش مهم‌ترین بخش برای روز دمو است و دقیق توضیح داده می‌شود.

\subsection{چطور برنامه واقعاً اجرا می‌شود}

هر اجرای هدف داخل یک کانتینر \lr{Docker} جداگانه انجام می‌شود. ابزار باینری و ورودی
را به مسیرهای ثابت \codeword|/binary| و \codeword|/case-input| سوار (\lr{mount}) می‌کند،
شبکه را قطع می‌کند، همهٔ \lr{capability}ها را می‌گیرد و محدودیت منابع اعمال می‌کند.
کد مربوط در \codeword|pipeline.py| (تابع \codeword|execute\_test|) به این شکل است:

\begin{latin}
\begin{lstlisting}[language=Python]
mode = test['input_mode']
input_value = '/case-input'
if mode == 'argv':          # ورودی به‌صورت آرگومان
    input_value = input_path.read_bytes().decode('utf-8')
command = [arg.replace('{binary}', '/binary').replace('{input}', input_value)
           for arg in config['input']['argv']]
box = Sandbox(config['environment']['image'], config['limits'])
result = box.execute(command, cwd,
                     stdin=input_path if mode == 'stdin' else None,
                     mounts=[(binary, '/binary'), (input_path, '/case-input')])
\end{lstlisting}
\end{latin}

محدودیت‌ها در پیکربندی مشخص می‌شوند (زمان دیواری، زمان پردازنده، حافظه، تعداد فرآیند و
حجم خروجی) و کانتینر با فلگ‌هایی مثل \codeword|--network none|، \codeword|--read-only|
و \codeword|--cap-drop ALL| ساخته می‌شود.

\subsection{تفکیک بیلد Release و Sanitized}

هر دو نسخه در دو حالت جداگانه ساخته و اجرا می‌شوند و نتایج آن‌ها مخلوط نمی‌شود:
\begin{itemize}
\item \textbf{Release}: با بهینه‌سازی و محافظ‌های کامپایلر
(\codeword|-O2 -D\_FORTIFY\_SOURCE=2 -fstack-protector-strong|) — نشان می‌دهد خطا در
باینری نهایی چطور بروز می‌کند.
\item \textbf{Sanitized}: با \codeword|-fsanitize=address,undefined| — جزئیات دقیق خطای
حافظه (نوع، محل، اندازهٔ دسترسی) را می‌دهد.
\end{itemize}

\subsection{طبقه‌بندی خطا}

ابزار چهار چیز را جدا نگه می‌دارد: \textbf{نوع خطا} (سرریز بافر هیپ/پشته، استفادهٔ پس
از آزادسازی و\ldots)، \textbf{روش تشخیص} (\lr{ASan}، \lr{UBSan}، محافظ قناری،
\lr{FORTIFY}، سیگنال سیستم‌عامل)، \textbf{نحوهٔ خاتمه} (کد خروج/سیگنال) و \textbf{محل}.
بخشی از کد مربوط که پیام \lr{ASan} را به نوع خطا تبدیل می‌کند:

\begin{latin}
\begin{lstlisting}[language=Python]
ASAN_TYPES = [
    ('heap-use-after-free', 'use-after-free'),
    ('heap-buffer-overflow', 'heap-buffer-overflow'),
    ('stack-buffer-overflow', 'stack-buffer-overflow'),
    ('SEGV', 'segmentation-fault'),
    ...
]
\end{lstlisting}
\end{latin}

نکتهٔ مهم: یک توقف بدون شواهد (مثلاً سیگنال خام بدون خروجی آشکارساز) هرگز به‌عنوان
آسیب‌پذیری تأییدشده ثبت نمی‌شود.

\subsection{مقایسه با مرجع مستقل}

صحت عملکرد با یک \lr{oracle} مستقل سنجیده می‌شود که از مستندات برنامه نوشته شده است
(نه از خروجی نسخهٔ وصله‌شده). برای هر آزمون، کد خروج، \lr{stdout} و \lr{stderr} مورد
انتظار به‌صورت بایت‌به‌بایت یا با الگو مشخص می‌شود:

\begin{latin}
\begin{lstlisting}[language=Python]
# oracles/libpng/poc-bit1-256pal.json  (خلاصه)
{ "exit_code": 1, "stdout": {"mode": "exact", "base64": ""},
  "stderr": {"mode": "exact", "base64": "cG5nIGVycm9yCg==" } }   # "png error\n"
\end{lstlisting}
\end{latin}

% ===================== پیوند یافته به وصله =====================
\section{اعتبارسنجی وصله و پیوند دادن یافته به تغییرات}

\subsection{تفکیک وضعیت یافته‌ها}

هر ورودی خطاساز روی هر دو نسخه اجرا می‌شود و یافته به یکی از این دسته‌ها تعلق می‌گیرد:
\begin{itemize}
\item \textbf{Fixed}: خطا روی پایه بازتولید می‌شود ولی روی وصله نه.
\item \textbf{Incomplete Fix}: خطا روی هر دو نسخه با همان امضا باقی می‌ماند (وصله ناقص).
\item \textbf{New Issue}: ورودی روی پایه پاک است ولی روی وصله خطای جدیدی می‌دهد.
\item \textbf{Functional Regression}: ورودی معتبر در پایه درست بود و روی وصله خراب شده.
\item \textbf{Inconclusive}: شواهد کافی نیست.
\end{itemize}

\subsection{همبستگی با diff}

صورت‌پروژه می‌گوید «محل خطا الزاماً محل تغییر نیست» و «هم‌جواری ایستا کافی نیست». در
نسخهٔ نهایی، ابزار همبستگی را سه‌سطحی می‌کند:
\begin{itemize}
\item \textbf{confirmed}: فقط وقتی خطِ \lr{crash} خودش خطِ تغییرکرده باشد (ارتباط
داده‌ای/کنترلی مستقیم).
\item \textbf{candidate}: تابع یا ماکروی تغییرکرده صرفاً روی پشتهٔ خطا حاضر است (نامزد،
نه تأیید).
\item \textbf{unknown}: شواهد کافی نیست.
\end{itemize}
وضعیت یافته (\lr{Fixed}/\lr{New Issue}) مستقل از وضعیت همبستگی می‌ماند. این دقیقاً همان
چیزی است که صورت‌پروژه می‌خواهد: «در نبود شواهد کافی باید \lr{unknown} گزارش شود».

\subsection{کدهای خروج مقرر}

تابع تصمیم (\codeword|decision.py|) تقدم کدهای خروج را پیاده می‌کند:

\begin{latin}
\begin{lstlisting}[language=Python]
def decide(axes):
    if axes.get('technical_error'):        return 30   # خطای ساخت/پیکربندی/ابزار
    if axes.get('confirmed_security_issue'): return 10   # VULN_REMAINS / NEW_ISSUE
    if axes.get('functional_regression'):   return 20   # پسرفت عملکردی
    if (axes.get('incomplete') or not axes.get('baseline_reproduced', False)
            or axes.get('hardening_degraded')): return 40   # شواهد ناکافی
    return 0                                             # PASS
\end{lstlisting}
\end{latin}

% ===================== تکرارپذیری و کوچک‌سازی =====================
\section{تکرارپذیری، کوچک‌سازی و آزمون رگرسیون}

\subsection{بستهٔ بازتولید}

به ازای هر خطای تأییدشده، ابزار یک بستهٔ بازتولید می‌سازد که شامل ورودی خطاساز،
شناسهٔ دقیق \lr{commit}، هش باینری، تنظیمات ساخت، متغیرهای محیطی و دستور اجراست.

\subsection{کوچک‌سازی خودکار ورودی}

ورودی خطاساز با الگوریتم \codeword|ddmin| (حذف تدریجی بایت‌ها) به‌صورت خودکار کوچک
می‌شود، به‌طوری که همان خطا با همان امضا حفظ شود. هم ورودی اولیه و هم ورودی کوچک‌شده
ذخیره می‌شوند.

\subsection{تولید خودکار آزمون رگرسیون}

ورودی کوچک‌شده به مجموعهٔ آزمون‌های رگرسیون
(\codeword|corpus/<target>/regressions/|) اضافه می‌شود و در اجراهای بعدی پایپ‌لاین
دوباره بررسی می‌گردد.

% ===================== سخت‌سازی =====================
\section{بررسی سخت‌سازی باینری}

ابزار وضعیت \lr{NX}، \lr{PIE}، \lr{RELRO} و \lr{Stack Canary} را با خروجی خام
\codeword|readelf| در بیلد \lr{Release} بررسی می‌کند و اگر وصله به‌طور ناخواسته یکی از
این محافظ‌ها را غیرفعال کند، هشدار می‌دهد (و این مانع صدور \lr{PASS} می‌شود). سطح
\lr{RELRO} به‌صورت بدون محافظت/جزئی/کامل گزارش می‌شود.

% ===================== اهداف =====================
\section{اهداف و سناریوهای پچ}

برای هر هدف چهار نسخه آماده شده است: \codeword|base| (آسیب‌پذیر)، \codeword|correct|
(وصلهٔ درست)، \codeword|incomplete| (وصلهٔ ناقص) و \codeword|regression| (وصلهٔ مخرب
عملکرد). علاوه بر این یک سناریوی \lr{New Issue} هم روی \codeword|tlvparse| اضافه شده
است.

\subsection{tlvparse — برنامهٔ خود گروه}

یک تحلیل‌گر قالب \lr{TLV}. آسیب‌پذیری آن سرریز بافر هیپ در \codeword|copy\_value|
است (کپی طول دلخواه در بافر ۳۲ بایتی).

\begin{center}
\begin{tabular}{lll}
\toprule
نسخه & تغییر & خروجی مورد انتظار \\
\midrule
\codeword|base| & بدون بررسی طول & آسیب‌پذیر \\
\codeword|correct| & بررسی طول در \codeword|copy\_value| & 0 (PASS) \\
\codeword|incomplete| & بررسی طول فقط در \codeword|load\_name| & 10 (VULN\_REMAINS) \\
\codeword|regression| & بررسی \codeword|>=| اشتباه که ورودی معتبر ۳۲ بایتی را رد می‌کند & 20 (REGRESSION) \\
\codeword|new-issue| & رفع \codeword|copy\_value| ولی افزودن بافر ۸ بایتی در \codeword|load\_blob| & 10 (NEW\_ISSUE) \\
\bottomrule
\end{tabular}
\end{center}

\subsection{giflib — پروژهٔ واقعی}

برنامهٔ \codeword|gif2rgb| از کتابخانهٔ واقعی \lr{giflib}. آسیب‌پذیری
\lr{CVE-2016-3977}: ایندکس رنگ زمینه بدون بررسی محدوده برای خواندن جدول رنگ استفاده
می‌شود (\lr{heap-buffer-overflow READ}). نسخهٔ \codeword|base| کامیت \codeword|c9a54af|
و \codeword|correct| کامیت \codeword|ea8dbc57| (وصلهٔ بالادستی) است.

\subsection{libpng — کتابخانهٔ واقعی با وصلهٔ ناقص واقعی}

این هدف به‌طور خاص سلسلهٔ «وصلهٔ ناقص سپس وصلهٔ کامل» را نشان می‌دهد:
\begin{itemize}
\item \lr{CVE-2015-8126}: \codeword|png\_set\_PLTE| اندازهٔ پالت را به عمق بیت محدود
نمی‌کرد.
\item \lr{CVE-2015-8472}: وصلهٔ ۱.۶.۱۹ ناقص بود (مسیر خواندن درست شد ولی
\codeword|png\_set\_PLTE| هنوز \codeword|png\_ptr| را به‌جای \codeword|info\_ptr| چک
می‌کرد) و در ۱.۶.۲۰ کامل شد.
\end{itemize}
نسخه‌ها: \codeword|base|=\lr{1.6.18}، \codeword|incomplete|=\lr{1.6.19}،
\codeword|correct|=\lr{1.6.20}. هر نسخه \textbf{درخت کامل بالادستی} است (نه زیرمجموعهٔ
دست‌چین‌شده).

\subsection{bytecount و lineparse — نمونه‌های بی‌خطر}

این دو فیچر برای خودآزمایی چارچوب هستند و همیشه خروجی 40 (شواهد ناکافی) می‌دهند.
نتیجهٔ آن‌ها \textbf{تأیید رفع یک آسیب‌پذیری واقعی نیست} و این نکته در \lr{README} هر
مخزن هم نوشته شده است.

% ===================== CI =====================
\section{اجرای خودکار در CI و وضعیت چک‌ها}

\subsection{مخزن اصلی: خودآزمایی}

فایل \codeword|.github/workflows/framework.yml| مخزن اصلی، کل ابزار را خودآزمایی
می‌کند: هر سناریوی پچ باید خروجی مورد انتظارش را بدهد (مثلاً \codeword|incomplete|
باید 10 بدهد). این خودآزمایی جدا از «پذیرش وصله» در مخازن هدف است.

\subsection{مخازن هدف: گیت پذیرش}

در هر مخزن هدف، پایپ‌لاین وصلهٔ \lr{PR} را ارزیابی می‌کند و چک فقط در صورتی سبز می‌شود
که اجرا کامل و معتبر باشد و خروجی 0 باشد. کدهای 10/20/30/40 چک را قرمز می‌کنند و نوع
نتیجه در پیام چک نمایش داده می‌شود. گزارش غایب/خراب یا کد ناشناخته هم هرگز سبز نمی‌شود.
نتیجهٔ واقعی (0/10/20/30/40) از سبز/قرمز بودن چک جدا است: چک سبز یعنی «ارزیابی درست
اجرا شد»، نه «وصله خوب است».

کامیت ابزار در مخازن هدف پین شده است تا نسخهٔ ارزیابی ثابت بماند.

% ===================== artifact =====================
\section{ساختار artifact}

هر اجرای پایپ‌لاین یک \lr{artifact} به نام \codeword|patchcheck-result| تولید می‌کند
که ساختارش این است:

\begin{latin}
\begin{lstlisting}[language=Bash]
artifacts/result/
  report.json        # گزارش ساختاریافتهٔ JSON (خروجی مقرر)
  summary.md         # گزارش خوانا برای انسان
  config.json        # تنظیمات موثر (نسخه‌ها، دستور ساخت، محدودیت‌ها)
  builds/            # هر ۴ بیلد (base/candidate × release/sanitized) + خروجی readelf
  builds.json        # هش باینری، فلگ‌ها، نتیجهٔ ساخت
  environment/       # پروب محیط (Linux x86_64، نسخهٔ gcc، ...)
  inputs/            # ورودی‌های آزمون و خطاساز
  runs/              # stdout/stderr/measurement هر اجرا
  replays/           # بازپخش‌ها (≥۳ بار برای خطاهای تأییدشده)
  sources/           # diff بین دو نسخه + درخت دقیق commitها
  bundles/           # بستهٔ بازتولید (reproduction.tar.gz)
\end{lstlisting}
\end{latin}

گزارش‌ها به فایل‌های داخل خود \lr{artifact} ارجاع می‌دهند، نه به مسیرهای بیرون از آن.
شواهد حتی در صورت شکست پایپ‌لاین هم ذخیره می‌شوند.

% ===================== راهنمای اجرا =====================
\section{راهنمای اجرای مستقل}

\begin{latin}
\begin{lstlisting}[language=Bash]
# 1) نصب وابستگی‌ها (آفلاین از wheelهای همراه)
python3 -m venv .venv && .venv/bin/python ci/bootstrap.py

# 2) ساخت مخازن fixture از پوشه targets/
python3 tools/setup_fixtures.py

# 3) ساخت تصویر ابزارزنجیره (Linux x86_64)
docker build --platform linux/amd64 -t patchcheck-framework:1.0 \
    -f environment/Dockerfile .

# 4) ارزیابی یک وصله
.venv/bin/patchcheck run configs/tlvparse-file.json --candidate incomplete --out artifacts/r

# 5) وصلهٔ نادیده (فایل diff)
.venv/bin/patchcheck run configs/giflib-file.json --patch /path/to/patch.diff --out artifacts/r

# 6) بازپخش / کوچک‌سازی / بسته‌بندی
.venv/bin/patchcheck replay  artifacts/r/report.json --test poc-long-blob --count 3
.venv/bin/patchcheck minimize artifacts/r/report.json --test poc-long-blob
.venv/bin/patchcheck bundle   artifacts/r/report.json --out reproduction.tar.gz
\end{lstlisting}
\end{latin}

% ===================== موارد دمو =====================
\section{موارد مهم برای روز دمو}

\begin{itemize}
\item \textbf{اجرای پویا}: نشان دهید که با خروجی \codeword|runs/| واقعاً باینری داخل
کانتینر اجرا می‌شود (نه تحلیل ایستا). خروجی \lr{ASan} در \codeword|runs/| برای بیلد
\codeword|sanitized| قابل نمایش است.
\item \textbf{چهار سناریو}: برای هر هدف، خروجی 0/10/20 برای
\codeword|correct|/\codeword|incomplete|/\codeword|regression| را نشان دهید و توضیح
دهید که چرا \codeword|incomplete| قرمز (10) است ولی یافته «ناقص» است.
\item \textbf{New Issue}: روی \codeword|tlvparse| شاخهٔ \codeword|new-issue| و ورودی
شاهد \codeword|poc-blob-new| را نشان دهید (پایه پاک، وصله خطای جدید).
\item \textbf{وصلهٔ نادیده}: با \codeword|--patch| یک \lr{diff} دلخواه را زنده ارزیابی
کنید.
\item \textbf{همبستگی}: توضیح دهید که \lr{confirmed} فقط برای خطِ \lr{crash} روی خطِ
تغییر است و بقیه \lr{candidate}/\lr{unknown} می‌ماند.
\item \textbf{سخت‌سازی}: خروجی \codeword|readelf| و مقایسهٔ NX/PIE/RELRO/canary را نشان
دهید.
\item \textbf{کوچک‌سازی و بازپخش}: اجرای \codeword|minimize| و \codeword|replay| را
نشان دهید.
\end{itemize}

% ===================== جمع‌بندی =====================
\section{جمع‌بندی}

ابزار \lr{Patchcheck} همهٔ بندهای صورت‌پروژه را پوشش می‌دهد: اجرای واقعی برنامهٔ هدف،
مقایسهٔ دو نسخه، پیوند یافته به تغییرات با شواهد واقعی، تفکیک چهار دسته یافته،
طبقه‌بندی خطا با تفکیک بیلد عادی/ابزارگذاری‌شده، تکرارپذیری و کوچک‌سازی، بررسی
سخت‌سازی باینری، و کدهای خروج استاندارد. ابزار روی پنج هدف مستقل آزمایش شده و
خروجی‌های واقعی آن‌ها با انتظار صورت‌پروژه مطابقت دارد.


\end{document}
